\documentclass[11pt,a4paper]{article}

\usepackage[utf8]{inputenc}
\usepackage{lmodern}
\usepackage[T1]{fontenc}
\usepackage[margin=2.5cm]{geometry}
\usepackage{amsmath,amssymb}
\usepackage{booktabs}
\usepackage{longtable}
\usepackage{tabularx}
\usepackage{microtype}
\usepackage[hidelinks]{hyperref}
\usepackage{xcolor}
\usepackage{listings}
\usepackage{authblk}

\lstset{
  basicstyle=\ttfamily\footnotesize,
  breaklines=true,
  frame=single,
  framesep=4pt,
  rulecolor=\color{gray!50},
  backgroundcolor=\color{gray!5},
  columns=fullflexible,
  keepspaces=true,
  language=Python,
  keywordstyle=\color{blue!60!black},
  commentstyle=\color{gray!70!black}\itshape,
  stringstyle=\color{green!40!black},
}

% Long \texttt identifiers must be allowed to break, or they overrun the
% measure in body text and in table cells.
\usepackage{seqsplit}
\newcommand{\code}[1]{\texttt{\seqsplit{#1}}}
\newcommand{\file}[1]{\texttt{\seqsplit{#1}}}

\emergencystretch=3em

\title{Code Narrative\\[0.4em]
\large An implementation guide to the agent-based model of\\
\emph{Driving in the wrong direction? Modelling policy mixes for EV adoption}}
\date{}
% Define authors with affiliation numbers
\author[1,2]{Miquel Bassart i Loré}
\author[3,4,5]{Daniel Torren-Peraire}

% Define affiliations matching the numbers
\affil[1]{Universität Bielefeld}
\affil[2]{Université Paris 1 Panthéon-Sorbonne}
\affil[3]{Universitat Autònoma de Barcelona}
\affil[4]{Università Ca' Foscari Venezia}
\affil[5]{International Institute for Applied Systems Analysis}


\begin{document}

\maketitle

\begin{abstract}
\noindent
This document is a guided tour of the code that implements the agent-based model
of the paper. It is organised to follow the model, one time step from beginning to end, then each
submodule in the order the paper introduces it and states throughout which
function implements which equation. Section~\ref{sec:notation} gives a symbol-to-variable
map, and Section~\ref{sec:extensions} documents mechanisms that exist in the code
but are switched off in every configuration used for the paper.
\end{abstract}

\bigskip
\hrule

\section{Orientation}

\subsection{Layout}

\begin{center}
\begin{tabularx}{\textwidth}{@{}lX@{}}
\toprule
Directory & Contents \\
\midrule
\file{package/model/} & The model itself: agents, landscapes, orchestration \\
\file{package/resources/} & Run harness (\code{generate\_data}, \code{load\_in\_controller}) and I/O helpers \\
\file{package/constants/} & \code{base\_params\_*.json} experiment configurations, \code{vary\_*.json} sweep definitions \\
\file{package/analysis/} & The paper's policy experiments (Section 4) \\
\file{package/generating\_data/} & Calibration, sensitivity and robustness runners \\
\file{package/plotting\_data/} & Figure scripts paired with the runners \\
\file{package/calibration/} & Empirical data assembly and simulation-based-inference calibration \\
\file{package/calibration\_data/} & California input series \\
\bottomrule
\end{tabularx}
\end{center}

\subsection{Running the model}

Dependencies are managed with \code{uv} (Python 3.13+); \code{uv sync} creates the
environment. Scripts resolve configuration paths relative to the repository root
and use absolute imports, so they are invoked as modules from the root:

\begin{lstlisting}[language=bash]
uv run python -m package.generating_data.calibration_gen
\end{lstlisting}

Each \code{*\_gen.py} script writes a timestamped directory under \code{results/}
and most call their matching \code{*\_plot.py} immediately afterwards. The
repository \code{README.md} carries the full map from paper figure to script and
configuration file.

\subsection{The single-run entry point}

Every experiment ultimately calls \code{generate\_data(parameters)} in
\file{package/resources/run.py}. It loads the pickled California input series,
derives the total horizon, constructs a \code{Controller}, and steps it to the end:

\begin{lstlisting}
parameters["time_steps_max"] = (parameters["duration_burn_in"]
                                + parameters["duration_calibration"]
                                + parameters["duration_future"])
controller = Controller(parameters)
while controller.t_controller < parameters["time_steps_max"]:
    controller.next_step()
return controller
\end{lstlisting}

The returned \code{Controller} \emph{is} the result: all recorded time series hang
off it (Section~\ref{sec:outputs}), and experiment scripts pickle it or extract
the series they need. \file{run.py} also holds a family of
\code{*\_parallel\_run} wrappers that map \code{generate\_data} over a list of
parameter dictionaries with \code{joblib}, each returning a different projection
of the results; these are what the experiment scripts actually call.

\section{One time step}
\label{sec:step}

\code{Controller.next\_step()} is short, and its order is the model's update
sequence:

\begin{lstlisting}
def next_step(self):
    self.update_time_series_data()
    self.cars_on_sale_all_firms = self.update_firms()
    self.second_hand_cars = self.get_second_hand_cars()
    self.consider_ev_vec, self.new_bought_vehicles = self.update_social_network()
    self.manage_saves()
    self.t_controller += 1
\end{lstlisting}

This maps one-to-one onto the numbered sequence in the paper's Figure~1:

\begin{center}
\begin{tabularx}{\textwidth}{@{}clX@{}}
\toprule
Step & Paper & Code \\
\midrule
0 & Exogenous variables and policies update & \code{update\_time\_series\_data()} \\
1 & Firms receive user and competitor information & \code{FirmManager.next\_step()} receives \code{consider\_ev\_vec} and \code{new\_bought\_vehicles} from the previous step \\
2 & Firms update product mix and prices & \code{Firm.choose\_cars\_segments()}, fired with probability $\theta$ \\
3 & Firms innovate & \code{Firm.innovate()}, fired with probability $\iota$ \\
4 & Used market scraps and reprices & \code{SecondHandMerchant.next\_step()} \\
5 & Users choose to keep, buy new, or buy used & \code{SocialNetworkUsers.update\_VehicleUsers()} \\
6 & Emissions are calculated & \code{update\_emisisons()}, within the user update \\
7 & Social parameters update for the next step & \code{calculate\_ev\_adoption()} \\
\bottomrule
\end{tabularx}
\end{center}

Two ordering details matter. First, firms act on information from the previous
period. \code{consider\_ev\_vec} and \code{new\_bought\_vehicles} are returned by
the user update at the \emph{end} of a step and consumed by the firm update at the
\emph{start} of the next, which makes the market-research channel a lag
rather than an instantaneous feedback. Second, EV consideration is recomputed
after all purchases are settled, so a user's threshold is evaluated against the
network state their own purchase has already contributed to.

\subsection{Exogenous inputs and policy application}

\code{update\_time\_series\_data()} indexes precomputed vectors at
\code{t\_controller} rather than recomputing paths per step. It sets the carbon
price, the effective gasoline cost, the grid emissions intensity, the electricity
price net of subsidy and gross of carbon price, the new and used EV rebates, and
the production subsidy. It also flips the EV research and production switches when
their start times arrive:

\begin{lstlisting}
if self.t_controller == self.ev_research_start_time:
    for firm in self.firm_manager.firms_list:
        firm.ev_research_bool = True
        firm.list_technology_memory = (firm.list_technology_memory_ICE
                                       + firm.list_technology_memory_EV)
\end{lstlisting}

The electricity price shows the two policy channels composing:
\begin{lstlisting}
self.electricity_price = (self.electricity_price_vec[t] * (1 - subsidy_prop)
                          + self.carbon_price * self.electricity_emissions_intensity)
\end{lstlisting}
The subsidy scales the pre-tax price; the carbon price is added on top of the
subsidised price, weighted by grid carbon intensity. Because intensity falls over
the policy period under the assumed linear decarbonisation, the carbon price
component of electricity shrinks as the grid cleans up, while the gasoline
component does not.

Policy values arrive as time series built by
\code{gen\_time\_series\_calibration\_scenarios\_policies()} from the
\code{parameters\_policies} block: \code{States} switches each instrument on or
off, \code{Values} sets its intensity. Instruments are applied from January 2024
at full intensity and held constant, so the series are step functions.

\section{Object model}

\code{Controller} owns the whole simulation and is the only object that knows the
calendar. It constructs and holds:

\begin{center}
\begin{tabularx}{\textwidth}{@{}lX@{}}
\toprule
Object & Responsibility \\
\midrule
\code{SocialNetworkUsers} & The whole user population as vectors, plus the network and the choice machinery \\
\code{FirmManager} & The firm population, market segmentation, and the competition statistics $W_{s,t}$ \\
\code{Firm} $\times J$ & One manufacturer: technology memory, product mix, pricing, innovation \\
\code{NKModel\_ICE}, \code{NKModel\_EV} & The two technology landscapes \\
\code{SecondHandMerchant} & The consolidated used car dealer and its stock \\
\bottomrule
\end{tabularx}
\end{center}

Note that users are \emph{not} individual objects in the hot path.
\code{SocialNetworkUsers} holds population attributes as NumPy vectors
(\code{beta\_vec}, \code{gamma\_vec}, \code{chi\_vec}, \code{d\_vec}) and evaluates
utility as a users-by-vehicles matrix; \code{VehicleUser} is a thin record used
for bookkeeping. Cars come in two kinds, and the distinction is load-bearing:
\code{CarModel} is a \emph{design} on sale (an attribute vector and a price, not a
physical car), while \code{PersonalCar} is an \emph{instance} that exists only once
a design is bought, since production is on demand. Both carry
\code{transportType}, which is \code{2} for ICEVs and \code{3} for EVs, a
convention worth memorising, as it appears in masks throughout the code.

\section{Car users}

\subsection{Choice set construction (Eq.~1)}

Two mechanisms restrict what a user may choose. Activation: with probability
$\eta$ (\code{prob\_switch\_car}) a user considers replacing their car at all;
otherwise their only option is what they already own. EV consideration: users
whose imitation threshold is not met cannot see EVs.

The second is implemented as a mask rather than by building per-user lists.
\code{gen\_mask()} constructs a users-by-vehicles boolean matrix by outer product,
and \code{masking\_options()} sets masked entries to $-\infty$ before
exponentiation:

\begin{lstlisting}
not_ev_vec = np.array([v.transportType == 2 for v in vehicles], dtype=bool)
consider_ev_all_vehicles = np.outer(consider_ev_vec == 1, ones)
dont_consider_evs_on_ev = np.outer(consider_ev_vec == 0, not_ev_vec)
ev_mask_matrix = (consider_ev_all_vehicles | dont_consider_evs_on_ev).astype(int)
\end{lstlisting}

A user who considers EVs gets every vehicle; a user who does not gets only ICEVs.
This is Eq.~1's case split, vectorised.

\subsection{Utility (Eq.~2--4)}

\code{vectorised\_calculate\_utility\_new\_cars()} and its second-hand counterpart
evaluate Eq.~2 as a matrix. The four terms are willingness to pay for quality with
diminishing returns $\alpha$, willingness to pay for range with diminishing
returns $\zeta$ (range being battery or tank size times efficiency, depreciated by
age), lifecycle emissions priced at $\gamma_i$, and lifecycle costs.

The discounted-sum terms of Eq.~3 and Eq.~4 share a single helper,
\code{\_lifecycle\_cost\_term()}, which implements the closed-form geometric series
of the naive-expectations case: a user assumes today's energy price and emissions
intensity persist indefinitely, so the infinite discounted sum collapses to
\[
D_i \frac{(1+r)(1-\delta_{a,t})\,c_{a,t}}{\omega_{a,t}\,(r - \delta_{a,t} - r\delta_{a,t})}.
\]
The purchase-price term is present only when the evaluated car is not the user's
own, and it is net of the expected sale price of the car being replaced.

Because a user's own car is always in the choice set, its utility is computed by a
separate path (\code{generate\_utilities\_current()}) with a small cache
(\code{\_build\_cv\_cache()}, \code{\_update\_cv\_cache\_row()}) so that owned-car
utility is not rebuilt from scratch each step.

\subsection{Logit choice (Eq.~5)}

Choice is sampled from the masked, exponentiated utilities. Users are activated
sequentially in random order, because a used car can only be bought once and the
stock must be updated between choices. \code{masking\_options()} subtracts the
row maximum before exponentiating for numerical stability, and
\code{user\_chooses()} samples from the unnormalised cumulative distribution:

\begin{lstlisting}
cumsum = np.cumsum(individual_specific_util_kappa)
choice_index = int(np.searchsorted(cumsum, u_draw * cumsum[-1], side='right'))
\end{lstlisting}

Scaling the uniform draw by \code{cumsum[-1]} avoids normalising the row, which is
where the $\kappa$-weighted denominator of Eq.~5 would otherwise be formed. If a
user has no admissible option, which happens only in the first steps, they keep
their current car.

The consequences of a choice are handled in the same function: the replaced car is
transferred to the used car dealer at its purchase price, unless it is an
initialisation car or already at scrap value, in which case it leaves the
simulation.

\subsection{Social imitation (Eq.~6--7)}

\code{create\_network()} builds a Watts--Strogatz small-world graph from
\code{SW\_network\_density} and \code{SW\_prob\_rewire}, stored as a normalised
adjacency matrix. \code{calculate\_ev\_adoption()} computes each user's
neighbourhood EV share and compares it against \code{chi\_vec}, returning the
binary \code{consider\_ev\_vec} that gates the next step's choice sets.

Thresholds are drawn from a beta distribution in \code{gen\_chi()}, with a fixed
\code{proportion\_zero\_target} (0.5\%) of the population pinned at $\chi_i = 0$ so
that contagion has a seed. Note that the threshold is a gate on
\emph{consideration}, not on choice: a user above their threshold still has to
prefer an EV on utility grounds, which is the behaviour--attitude gap the paper
discusses.

\section{Car supply}

\subsection{Segmentation (Appendix A.2.1)}

\code{FirmManager} partitions the population into segments on willingness to pay
for quality (\code{num\_beta\_segments}, $R^\beta = 8$) and for emissions reduction
(\code{num\_gamma\_segments}, $R^\gamma = 2$), crossed with the binary EV
consideration state. Each segment is represented by preference parameters at the
corresponding quantile, giving the representative consumer profiles that firms
price against.

\subsection{Competition state (Eq.~9)}

\code{update\_W\_immediate()} computes the current $\kappa$-weighted sum of
lifetime utilities on offer per segment, and
\code{update\_market\_data\_moving\_average()} maintains the twelve-month moving
average that becomes $W_{s,t}$. Crucially, for segments that do not consider EVs
only ICE cars enter the sum, matching Eq.~9's case split. The result is passed
down to firms as \code{W\_vec} together with \code{I\_s\_t\_vec}, the segment
populations.

\subsection{Optimal pricing (Eq.~10)}

\code{Firm.calc\_optimal\_price\_cars()} implements Eq.~10 directly, including its
Lambert~$\mathbb{W}$ term, as a cars-by-segments matrix operation:

\begin{lstlisting}
U = term1 + term2 + term3 + term4          # utility at P = marginal cost
exp_input = (self.kappa * U - 1) - np.log(self.W_vec[np.newaxis, :])
LW = lambertw(np.exp(exp_input), 0).real
P = C_m_cost[:, np.newaxis] + (1.0 + LW) / self.kappa
\end{lstlisting}

Two subsidies enter here, and they enter differently. This is exactly the
distributional distinction the paper draws between a production subsidy and a
purchase rebate:

\begin{lstlisting}
C_m_cost[ev_mask]  = np.maximum(0, C_m[ev_mask] - self.production_subsidy)
C_m_price[ev_mask] = np.maximum(0, C_m[ev_mask]
                                - (self.production_subsidy + self.rebate
                                   + self.rebate_calibration))
\end{lstlisting}

\code{C\_m\_cost} is the cost the firm actually bears, and it sets the price floor.
\code{C\_m\_price} is the cost relevant to the consumer's perceived price, and it
enters the utility term. The production subsidy appears in both, so the firm
captures part of it as rent; the rebate appears only in the second, so it passes
through to the buyer. Both are floored at zero, which is the zero-lower-bound the
paper describes and the mechanism behind the finding that stacking a rebate on a
production subsidy only inflates profits once the bound binds.

\subsection{Product mix (Appendix A.2.2)}

\code{choose\_cars\_segments()} implements the sequential heuristic: build the
segments-by-technologies expected-profit matrix
(\code{set\_utility\_and\_profit\_grid\_production()},
\code{calc\_predicted\_profit\_segments\_production()}), repeatedly take the most
profitable cell, add that technology at that segment's optimal price, delete the
segment row, lock the price for that technology across remaining segments, and
recompute, repeating until \code{max\_cars\_prod} ($Y^+$) is reached or the matrix
empties.
Locking the price is what allows a design to serve several segments at one price
while being less competitive in the poorer ones.

The whole procedure fires only with probability $\theta$
(\code{prob\_change\_production}), representing price rigidity and information
delay.

\subsection{Innovation (Eq.~11--15)}

\code{innovate()} fires with probability $\iota$ (\code{prob\_innovate}), which is
how failed innovation is represented. It builds the candidate set of Eq.~11 from
unexplored technologies at Hamming distance one from the current search location
on each landscape, plus the current locations themselves so that not searching
remains an option (\code{generate\_neighbouring\_technologies()},
\code{gen\_neighbour\_strings()}). Candidates are scored by expected profit at
their best segment (\code{calc\_predicted\_profit\_segments\_research()}) and one
is drawn by a logit with research intensity $\lambda$
(\code{select\_car\_lambda\_research()}), balancing exploration against
exploitation.

Memory follows Eq.~14--15. \code{update\_memory\_timer()} increments a per-technology
idle counter, zeroed for anything researched or in production, and
\code{update\_memory\_len()} evicts the longest-idle technology not in production
once \code{memory\_cap} ($M^+$) is exceeded.

\subsection{NK landscapes (Supplementary 3.1)}

\file{nkModel\_ICE.py} and \file{nkModel\_EV.py} are near-identical, differing in
dimensionality: both map an $N$-bit string with epistasis $K$ to attributes, but
the EV landscape carries battery size as an additional attribute, so ICE has three
attributes and EV four. \code{generate\_fitness\_landscape()} builds the
contribution tables, and \code{calculate\_fitness\_vectorized()} evaluates
attributes for many strings at once. Attribute ranges and the cost-attribute
correlations $\rho$ come from the configuration.

Landscapes are constructed once per run by
\code{setup\_ICE\_landscape()}/\code{setup\_EV\_landscape()} and shared across
firms, so firms differ only in where they have searched, which is the source of
emergent specialisation. \code{find\_min\_fitness\_string()} provides the
initialisation point: firms start one Hamming step from the worst technology
evaluated at the median segment.

\section{Used car market (Eq.~16--21)}

\code{SecondHandMerchant} is the single consolidated dealer. Its pricing heuristic
(\code{calc\_car\_price\_heuristic()}) implements Eq.~18--20: normalise quality,
efficiency and battery size by the maxima over new cars, find each used car's
nearest new design in that space, take its price net of rebates, and depreciate by
age at the type-specific monthly rate \code{delta\_P} ($\xi$):

\begin{lstlisting}
closest_prices = np.maximum(first_hand_prices[closest_idxs]
                            - (self.rebate_calibration + self.rebate), 0)
adjusted_prices = closest_prices * (1 - second_hand_delta_P) ** second_hand_ages
\end{lstlisting}

Because EVs depreciate faster than ICEVs ($\xi$ of 0.87\% against 1.16\% monthly),
the used EV stock holds value differently, which is what makes the used EV rebate
weak early in the transition: there is very little used EV stock to subsidise.

\code{update\_stock\_contents()} applies the markup $\mu$ to form purchase prices
(Eq.~16), scraps cars priced below scrap value $F$ (Eq.~21), and enforces the two
computational bounds: \code{max\_num\_cars\_prop} ($V^+$, the stock ceiling, 10\% of
population) and \code{age\_limit\_second\_hand} ($T^+$, months unsold before
scrapping). The dealer also refreshes owned-car fuel costs from the current
gas price each step, which is why a driving-cost shock reaches second-hand cars
too.

\section{Run structure and controller reuse}
\label{sec:phases}

A run has three phases, measured in months and set by configuration:

\begin{center}
\begin{tabularx}{\textwidth}{@{}llX@{}}
\toprule
Phase & Length & Purpose \\
\midrule
\code{duration\_burn\_in} & 180 & ICEV-only market reaches a steady state; EV research and production are disabled \\
\code{duration\_calibration} & 276 & 2001--2023, matched against California data \\
\code{duration\_future} & 144 & 2024--2035 policy period, extended to 2050 for the stability analysis \\
\bottomrule
\end{tabularx}
\end{center}

\code{manage\_burn\_in()}, \code{manage\_calibration()} and \code{manage\_scenario()}
switch the exogenous regime between phases. Burn-in uses \code{next\_step\_burn\_in}
variants on the firm side, which update product mix without the EV machinery.

The important performance property is that the first two phases are computed once
and reused. \code{setup\_continued\_run\_future()} takes an already-calibrated
controller, applies a new future-period policy configuration, and
\code{load\_in\_controller()} continues stepping from where calibration ended:

\begin{lstlisting}
def load_in_controller(controller_load, base_params_future):
    base_params_future["time_steps_max"] = (controller_load.duration_burn_in
                                           + controller_load.duration_calibration
                                           + base_params_future["duration_future"])
    controller_load.setup_continued_run_future(base_params_future)
    while controller_load.t_controller < base_params_future["time_steps_max"]:
        controller_load.next_step()
    return controller_load
\end{lstlisting}

Policy scenarios are therefore deep copies of one calibrated controller, which is
what makes a Bayesian optimisation over policy intensity affordable: 456 months of
burn-in and calibration are not re-simulated per candidate intensity. It also
means every policy scenario shares an identical history up to December 2023, so
differences between scenarios are attributable to policy rather than to
stochastic divergence before it.

\section{Policy instruments}

The five instruments of the paper enter at different points, summarised here
because the entry point determines the economics:

\begin{center}
\begin{tabularx}{\textwidth}{@{}lX@{}}
\toprule
Instrument & Where it acts \\
\midrule
Carbon price & Added to gasoline and electricity cost per kWh, weighted by emissions intensity, in \code{update\_time\_series\_data()}. Affects all users every month. \\
Electricity subsidy & Scales the pre-tax electricity price. Affects EV users every month. \\
New EV rebate & Reduces the consumer-facing cost term \code{C\_m\_price} in pricing; floored at zero. One-off at purchase. \\
Used EV rebate & Reduces used EV prices in the dealer's heuristic and in second-hand utility. One-off at purchase. \\
Production subsidy & Reduces both \code{C\_m\_cost} and \code{C\_m\_price}, so firms retain part of it. One-off at purchase. \\
\bottomrule
\end{tabularx}
\end{center}

Net public cost is accumulated as \emph{policy distortion} on each actor and
aggregated by \code{calc\_total\_policy\_distortion()} and
\code{calc\_net\_policy\_distortion()}, the latter negated so that a positive value
is a cost to the public purse. Carbon price revenue and subsidy outlays therefore
net against each other, which is how the near-budget-neutral carbon-price-plus-
electricity-subsidy combination arises.

Note that the calibration-period rebates, the simplified federal and state EV
incentives of 2010--2023, are carried by separate
\code{rebate\_calibration}/\code{used\_rebate\_calibration} series so that they can
be switched off independently of the policy-period instruments, as the paper's
experimental design requires.

\section{Calibration}

Calibration proceeds in the layers the paper describes.

\code{calibration\_data\_inputs.py} assembles the California series (vehicle
population, EV sales, gasoline and electricity prices, grid emissions intensity,
CPI) into the pickled object \code{generate\_data} loads;
\code{calibration\_data\_outputs.py} formats the observed targets.
\code{fit\_distance.py} fits the monthly distance-driven distribution.

Analytically determined parameters are computed inside the controller rather than
configured. \code{calc\_beta\_median()} solves for the median willingness to pay
for quality such that the optimal price of an average ICEV targeted at an average
consumer equals the observed average price; \code{generate\_beta\_values\_quintiles()}
then distributes it across the population by the Californian income
distribution, so richer agents get higher $\beta_i$. The remaining population
draws are \code{gen\_distance()}, \code{gen\_chi()}, \code{gen\_nu()} and
\code{gen\_gamma()}.

The innovativeness threshold distribution is calibrated by simulation-based
inference in \code{NN\_multi\_round\_calibration\_multi\_gen.py}, which uses
\code{sbi} to fit the beta distribution parameters \code{a\_chi} and \code{b\_chi}
against 2010--2023 EV uptake over multiple rounds. Consumer intensity of choice
$\kappa$ and the efficiency depreciation rate $\delta$ are set by grid search
against the target ranges for price, market concentration and car age.

\section{Outputs}
\label{sec:outputs}

Recording is opt-in and thinned. \code{save\_timeseries\_data\_state} enables it and
\code{compression\_factor\_state} sets the stride, so long multi-seed runs can
record sparsely or not at all. Each object has a
\code{set\_up\_time\_series\_*}/\code{save\_timeseries\_*} pair, and
\code{Controller.manage\_saves()} drives them.

Series live on the object that owns the quantity: adoption, prices, emissions,
utility and car age on \code{social\_network}; concentration, profit and margins on
\code{firm\_manager}. The scalar summaries used as optimisation targets and table
entries are computed on demand (\code{calc\_EV\_prop()},
\code{calc\_mean\_car\_age()}, \code{calc\_last\_step\_HHI()},
\code{calc\_price\_range\_ice()}), and the aggregation helpers in \file{run.py}
select which of these to return across a parameter sweep.

\section{Extensions present in the code but not used in the paper}
\label{sec:extensions}

The model contains three optional mechanisms that are disabled by default and
switched off in every configuration in this repository. They produce no result
reported in the paper. Each is guarded so that omitting its parameter reproduces
the published behaviour exactly; they are retained as a foundation for follow-up
work on command-and-control policy and on policy anticipation.

\subsection{ICE bans}

The paper restricts itself to market-based instruments and explicitly excludes
command-and-control policy. Three independent ban levers are nevertheless
implemented, each expressed as months after the end of burn-in, the same
convention as \code{ev\_production\_start\_time}, and each defaulting to
\code{None}, meaning never:

\begin{center}
\begin{tabularx}{\textwidth}{@{}lX@{}}
\toprule
Parameter & Effect \\
\midrule
\code{ICE\_research\_ban\_time} & Firms may no longer research or improve ICE technology. \code{innovate()} stops generating ICE candidates. Existing ICE designs stay in memory and remain sellable. \\
\code{ICE\_sales\_ban\_time} & Firms may no longer offer new ICEVs. \code{choose\_cars\_segments()} drops ICE from the candidate pool, and \code{next\_step()} immediately flushes any ICE design left on sale rather than waiting for the next $\theta$ roll. Already-sold ICEVs remain drivable and resellable. \\
\code{ICE\_driving\_ban\_time}, \code{ICE\_driving\_ban\_penalty} & A per-unit penalty added to effective ICE fuel cost from the ban date onward, folded into \code{gas\_cost\_effective\_vec}. Reaches second-hand ICEVs, since the dealer refreshes fuel costs from the same series. \\
\bottomrule
\end{tabularx}
\end{center}

The three are additive, so research-only, research-plus-sales, and
research-plus-sales-plus-driving scenarios can all be configured. The design
choice worth noting is that the driving ban is a cost shock rather than a
prohibition: it therefore flows through the existing utility-driven choice
mechanism and requires no new decision rule on either side of the market, whereas
the sales and research bans are hard constraints on firms' choice sets. All three
validate that EV research or production has already begun before the ban binds.
The active flags are recomputed every step from the configured times rather than
latched by a one-shot equality test, so they survive
\code{setup\_continued\_run\_future()} re-deriving them for a new scenario.

\subsection{Forward-looking expectations}

\code{forward\_looking\_expectations} (default \code{False}) switches agents away
from the naive expectations of the paper, under which the current energy price and
emissions intensity are assumed to persist and the discounted sum collapses to the
closed form of Section~\ref{sec:step}. With the flag on, agents instead discount
the actual known future path.

\code{compute\_discounted\_indices()} precomputes present-value indices by backward
recursion, $\mathrm{Index}(t) = z_t + q\,\mathrm{Index}(t+1)$, once per run rather
than per agent, so the switch costs nothing per utility evaluation. Unrolling the
recursion with a constant $z_t$ reproduces the naive closed form exactly, so it is
a strict generalisation. With the flag off the indices are still computed but
never read by the utility formulas.

The boundary condition is the subtle part. A flat-hold extension past the
simulated horizon would be wrong for the carbon price specifically, because
\code{calculate\_price\_at\_time()} returns zero once the policy period ends: the
schedule is coded to stop, not to persist. Since \code{duration\_future} typically
ends at the policy horizon, a flat hold would have a forward-looking agent treat a
temporary tax as permanent at its peak rate. The carbon component is therefore
extended using its own schedule, while base energy prices and the decarbonisation
trend are extended by holding their final value, which for those series is the
correct continuation rather than an approximation.

This is the channel through which forward-looking agents would anticipate an
announced ban or a scheduled carbon price before it arrives, over a horizon
implied by the model's own discount and depreciation rates rather than by a
separate anticipation parameter.

\subsection{Carbon price schedules and the research subsidy}

\code{calculate\_growth()} supports \code{flat}, \code{linear}, \code{quadratic}
and \code{exponential} carbon price paths via \code{Carbon\_price\_state}. Every
configuration here uses \code{flat}, matching the paper's constant-intensity
design, but the other schedules are the hook for time-varying schemes such as the
EU~ETS. A sixth policy lever, \code{Research\_subsidy}, is present in the policy
state dictionary but is not among the five instruments analysed.

\section{Implementation notes}

\paragraph{Vectorisation.} The hot path is the users-by-vehicles utility matrix,
rebuilt each step. Utility, masking, exponentiation and pricing are all matrix
operations; the only per-agent Python loop is the sequential choice itself, which
cannot be vectorised because used car stock must update between choices.

\paragraph{Randomness and reproducibility.} Two seeds are separated deliberately.
\code{seed\_inputs} governs structural draws (landscapes, network, population
attributes), and \code{seed} governs run-time stochasticity: activation,
innovation and production rolls, and choice draws. \code{handle\_seed()}
distributes \code{RandomState} instances to each component. Fixing
\code{seed\_inputs} while varying \code{seed} therefore isolates behavioural noise
from structural variation, which is what multi-seed experiments exploit.

\paragraph{Parallelism.} Monte Carlo replication is parallel across seeds via
\code{joblib}, with worker count from \code{get\_num\_workers()}. Everything is
CPU-bound; \code{torch} is present only because \code{sbi} requires it for the
inference-based calibration, and is pinned to a CPU-only build.

\paragraph{Cost.} A single run is 600 months of a 3{,}000-user, 16-firm market with
two NK landscapes. The policy experiments dominate total cost, since each Bayesian
optimisation iteration is 64 seeded runs of the future period; controller reuse
(Section~\ref{sec:phases}) is what keeps this tractable.

\section{Symbol-to-variable map}
\label{sec:notation}

\begin{longtable}{@{}p{0.13\textwidth}p{0.47\textwidth}p{0.32\textwidth}@{}}
\toprule
Symbol & Code & Location \\
\midrule
\endfirsthead
\toprule
Symbol & Code & Location \\
\midrule
\endhead
$\beta_i$ & \code{beta\_vec}, \code{beta\_s\_values} & social network, firm segments \\
$\gamma_i$ & \code{gamma\_vec}, \code{gamma\_s\_values} & social network, firm segments \\
$\nu$ & \code{nu} & vehicle user parameters \\
$\alpha$, $\zeta$ & \code{alpha}, \code{zeta} & vehicle user parameters \\
$\kappa$ & \code{kappa} & vehicle user parameters \\
$\chi_i$ & \code{chi\_vec}; \code{a\_chi}, \code{b\_chi} & social network \\
$D_i$ & \code{d\_vec}; \code{d\_mean} & social network; firm \\
$r$ & \code{r} & vehicle user parameters \\
$\delta$ & \code{delta} & ICE/EV parameters \\
$\xi$ & \code{delta\_P} & ICE/EV parameters \\
$\eta$ & \code{prob\_switch\_car} & social network \\
$\iota$ & \code{prob\_innovate} & firm \\
$\theta$ & \code{prob\_change\_production} & firm \\
$\lambda$ & \code{lambda} & firm \\
$\mu$ & \code{mu} & second-hand market \\
$M^+$ & \code{memory\_cap} & firm \\
$Y^+$ & \code{max\_cars\_prod} & firm \\
$V^+$ & \code{max\_num\_cars\_prop} & second-hand market \\
$T^+$ & \code{age\_limit\_second\_hand} & second-hand market \\
$F$ & \code{scrap\_price} & second-hand market \\
$N$, $K$ & \code{N}, \code{K} & ICE/EV landscape \\
$E$ & \code{production\_emissions} & ICE/EV parameters \\
$J$, $I$ & \code{J}, \code{num\_individuals} & firm manager, social network \\
$R^\beta$, $R^\gamma$ & \code{num\_beta\_segments}, \code{num\_gamma\_segments} & firm manager \\
$\rho_{WS}$, $p_{WS}$ & \code{SW\_network\_density}, \code{SW\_prob\_rewire} & social network \\
$Q_{a,t}$, $\omega_{a,t}$, $B_{a,t}$ & \code{Quality\_a\_t}, \code{Eff\_omega\_a\_t}, \code{B} & car attribute dicts \\
$L_{a,t}$ & \code{L\_a\_t} & car objects \\
$z_{v,t}$ & \code{transportType} (2 = ICE, 3 = EV) & car objects \\
$\upsilon_{v,t}$ & \code{owner\_id} & car objects \\
$x_{i,t}$ & \code{consider\_ev\_vec} & social network \\
$\psi_{i,t}$ & \code{vehicle\_chosen} & \code{user\_chooses()} \\
$U_{a,i,t}$ & \code{utilities\_matrix} & social network \\
$W_{s,t}$ & \code{W\_vec}, \code{W\_segment} & firm manager \\
$I_{s,t}$ & \code{I\_s\_t\_vec} & firm manager \\
$P^*_{s,m,t}$ & \code{optimal\_price\_segments} & car objects \\
$M_{j,t}$ & \code{list\_technology\_memory} & firm \\
$M^{RD}_{j,t}$ & \code{generate\_neighbouring\_technologies()} & firm \\
$T_{m,j,t}$ & \code{update\_memory\_timer()} & firm \\
\bottomrule
\end{longtable}

\end{document}
